
\setcounter{footnote}{29}

take \(B=B'\) to be the Borel subgroup defined by \(c=0\), and let \(u\)
be the following automorphism of \(B\):
\[
  u(a,b,d,h)=(a,b,d,ah).
\]

\par\medskip\noindent\SGAInlineRefTarget{sga3:xxiv:corollary:6-4}\textbf{Corollary 6.4.---}\enspace
\emph{The functor
\[
  (G,B)\longmapsto B
\]
from the category of pairs \((G,B)\), where \(G\) is an
\(S\)-semisimple group and \(B\) is a Borel subgroup of \(G\), to the
category of \(S\)-group schemes, with isomorphisms as morphisms, is
fully faithful.}

\par\medskip\noindent\SGAInlineRefTarget{sga3:xxiv:corollary:6-5}\textbf{Corollary 6.5.---}\enspace
\emph{Let \(S\) be a scheme, \(G\) an \(S\)-semisimple group, \(B\) a
Borel subgroup of \(G\), and \(B'\) an \(S\)-group locally isomorphic to
\(B\) for the \textup{(fpqc)} topology. Then \(B'\) is locally
isomorphic to \(B\) for the local finite étale
topology,\footnote{\textit{Editors' note.} This topology is denoted
\textup{(étf)} in \SGARefLink{sga3:iv:subsection:6:3}{Exposé IV, 6.3}. In other words, for every \(s\in S\)
there exist an open neighborhood \(U\) of \(s\) and a finite surjective
étale morphism \(V\to U\) such that
\(B'\times_S V\simeq B\times_S V\).} and there exists an
\(S\)-semisimple group \(G'\) of which \(B'\) is a Borel subgroup.
Moreover, \(G'\) is unique up to a unique isomorphism inducing the
identity on \(B'\).}

\par\medskip\noindent\SGAInlineRefTarget{sga3:xxiv:corollary:6-6}\textbf{Corollary 6.6.---}\enspace
\emph{Let \(S\) be a scheme, \(G\) an \(S\)-semisimple group, and \(B\)
a Borel subgroup of \(G\). Then
\(\underline{\operatorname{Aut}}_{S\text{-}\mathrm{gr.}}(B)\) is
representable by an affine smooth \(S\)-scheme,
\(\underline{\operatorname{Autext}}(B)\) is representable by an étale
finite \(S\)-scheme, and the evident morphisms induce an isomorphism of
exact sequences.}

\[
  \begin{tikzcd}[column sep=small,row sep=large]
1 \arrow[r]
&
\underline{B}_{\mathrm{ad}}(G)
  \arrow[r]
  \arrow[d,"\wr" description]
&
\underline{\Aut}_{S\text{-}\mathrm{gr.}}(G,B)
  \arrow[r]
  \arrow[d,"\wr" description]
&
\underline{\Autext}(G)
  \arrow[r]
  \arrow[d,"\wr" description]
&
1
\\
1 \arrow[r]
&
\underline{B}_{\mathrm{ad}}
  \arrow[r]
&
\underline{\Aut}_{S\text{-}\mathrm{gr.}}(G)
  \arrow[r]
&
\underline{\Autext}(B)
  \arrow[r]
&
1 .
\end{tikzcd}

\]

This follows immediately from \SGARefLink{sga3:xxiv:2.1}{2.1} and the preceding results. We leave
to the reader the task of developing the analogues of the results of
nos.~1, 2, 3, and 4 in the foregoing setting.

\par\medskip\noindent\SGAInlineRefTarget{sga3:xxiv:remark:6-7}\textbf{Remark 6.7.---}\enspace
If \(S\) is a scheme and \(B\) an \(S\)-group, then \(B\) can be a
Borel subgroup of only one semisimple group, uniquely determined by
\(B\). It remains to characterize those \(S\)-groups \(B\) which are
Borel subgroups of semisimple
groups.\footnote{\textit{Editors' note.} One may ask whether every
smooth affine \(S\)-group whose geometric fibers are Borel subgroups of
semisimple groups is itself a Borel subgroup of a semisimple
\(S\)-group. We have suppressed the counterexample in the original,
with \(S\) the scheme of dual numbers over a field \(k\), because it
was erroneous, as M.~Demazure pointed out.}

\section*{\SGAInlineRefTarget{sga3:xxiv:section:7}7. Representability of the functors
\(\underline{\operatorname{Hom}}_{S\text{-}\mathrm{gr.}}(G,H)\), for
\(G\) reductive}

\subsection*{\SGAInlineRefTarget{sga3:xxiv:subsection:7-1}7.1. The split case}

\subsubsection*{7.1.1.}

Let \(S\) be a scheme, \(G\) a pinned \(S\)-group, \(T\) its maximal
torus, and \(\Delta\) the set of simple roots. For
\(\alpha\in\Delta\), let
\[
  u_\alpha\in U_\alpha^\times(S),
  \qquad
  w_\alpha\in\operatorname{Norm}_G(T)(S)
\]
be the elements defined by the pinning.

Let \(H\) be an \(S\)-group scheme. We are interested in the functor
\(\underline{\operatorname{Hom}}_{S\text{-}\mathrm{gr.}}(G,H)\), and
more precisely in the morphism
\[
  q:\underline{\operatorname{Hom}}_{S\text{-}\mathrm{gr.}}(G,H)
  \longrightarrow
  \underline{\operatorname{Hom}}_{S\text{-}\mathrm{gr.}}(T,H).
\]
Let
\[
  f_T:T\longrightarrow H
\]
be an \(S\)-group homomorphism and put \(q^{-1}(f_T)=\mathscr F\).
Thus
\[
  \mathscr F(S')
  =
  \left\{
    f\in\operatorname{Hom}_{S'\text{-}\mathrm{gr.}}
      (G_{S'},H_{S'})
    \,\middle|\,
    f|_{T_{S'}}=(f_T)_{S'}
  \right\}.
\]
There is a morphism of \(S\)-functors
\[
  i:\mathscr F\longrightarrow H^{2n},
\]
where \(n=\operatorname{Card}(\Delta)\), defined on points by
\[
  i(f)=\bigl(f(u_\alpha),f(w_\alpha)\bigr)_{\alpha\in\Delta}.
\]
By \SGARefLink{sga3:XXIII:1.9}{Exposé XXIII, 1.9}, \(i\) is a monomorphism.

\par\medskip\noindent\SGAInlineRefTarget{sga3:xxiv:proposition:7-1-2}\textbf{Proposition 7.1.2.---}\enspace
\emph{If \(H\) is separated over \(S\), then \(\mathscr F\) is
representable and \(i\) is a closed immersion.}

The usual technique of relative
representability\footnote{\textit{Editors' note.} See, for example, the
proof of \SGARefLink{sga3:xxii:thm:5.8.1}{Exposé XXII, 5.8.1}.} shows that it is enough to prove the
following. Given sections
\[
  v_\alpha,h_\alpha\in H(S),
  \qquad \alpha\in\Delta,
\]
the \(S\)-schemes \(S'\) for which there exists an \(S'\)-homomorphism
\[
  f:G_{S'}\longrightarrow H_{S'}
\]
extending \((f_T)_{S'}\) and satisfying
\[
  f(u_\alpha)=v_\alpha,
  \qquad
  f(w_\alpha)=h_\alpha
\]
are exactly those which factor through a certain closed subscheme of
\(S\). We may plainly suppose \(S\) affine.

For convenience, say that a morphism \(Y\to
X\)\footnote{\textit{Editors' note.} We have corrected \(X\to Y\) to
\(Y\to X\).} of affine schemes satisfies condition \textup{(L)} if
\(Y\) is the spectrum of an \(\mathcal O(X)\)-algebra which is a free
\(\mathcal O(X)\)-module. Within the category of affine schemes,
products and composites of morphisms satisfying \textup{(L)} again
satisfy \textup{(L)}, and \textup{(L)} is stable under base change.

\par\medskip\noindent\SGAInlineRefTarget{sga3:xxiv:lemma:7-1-3}\textbf{Lemma 7.1.3.---}\enspace
\emph{Suppose that \(S\) is affine, and let \(\alpha\in\Delta\).
Consider the morphism
\[
  a:T\times_S T\longrightarrow G_{a,S},
  \qquad
  a(t,t')=\alpha(t)+\alpha(t').
\]
Then:
\begin{enumerate}
\item[(i)] \(a\) is faithfully flat and of finite presentation;
\item[(ii)] if \(R\) is the fiber square of \(a\), the structure
  morphism \(R\to S\) satisfies \textup{(L)}.
\end{enumerate}}

The morphism \(a\) is surjective, since
\(\alpha:T\to G_{m,S}\) is surjective, and it is plainly of finite
presentation. Since \(\alpha\) satisfies \textup{(L)}, it is enough,
for both assertions, to show that
\[
  u:G_{m,S}^2\longrightarrow G_{a,S},
  \qquad
  u(x,y)=x+y
\]
satisfies \textup{(L)}. Equivalently, for every ring \(A\), the ring
\[
  A[X,Y,X^{-1},Y^{-1}]
\]
must be free over its subring \(A[X+Y]\).\footnote{\textit{Editors'
note.} We have simplified the proof in the original.} Now \(A[X,Y]\)
is free with basis \(\{1,X\}\) over
\[
  B=A[X+Y,XY],
\]
because \(X^2-(X+Y)X+XY=0\). Hence
\[
  A[X^{\pm1},Y^{\pm1}]
  =A[X,Y]_{XY}
\]
is free over
\[
  B_{XY}=A[X+Y,(XY)^{\pm1}]
\]
with basis \(\{1,X\}\), and free over \(A[X+Y]\) with basis
\[
  (XY)^p,\quad X(XY)^p,
  \qquad p\in\mathbb Z.
\]

\par\medskip\noindent\SGAInlineRefTarget{sga3:xxiv:lemma:7-1-4}\textbf{Lemma 7.1.4.---}\enspace
\emph{Let \(\alpha\in\Delta\), and let
\[
  b:T\times_S T\longrightarrow H
\]
be the morphism defined on points by
\[
  b(t,t')
  =
  \operatorname{int}(f_T(t))(v_\alpha)\,
  \operatorname{int}(f_T(t'))(v_\alpha).
\]
Let \(f_\alpha:U_\alpha\to H\) be an \(S\)-morphism. The following
conditions are equivalent:\footnote{\textit{Editors' note.} In what
follows the group law of \(U_\alpha\) is written additively. If
\(p_\alpha:G_{a,S}\xrightarrow{\sim}U_\alpha\) is the isomorphism with
\(p_\alpha(1)=u_\alpha\), and \(x=p_\alpha(z)\), then
\(\alpha(t)x\) denotes \(p_\alpha(\alpha(t)z)\), while
\(a(t,t')u_\alpha\) denotes
\(p_\alpha(a(t,t'))=\alpha(t)u_\alpha+\alpha(t')u_\alpha\).}
\begin{enumerate}
\item[\SGAInlineRefTarget{sga3:xxiv:lemma:7-1-4:item:i}(i)] \(f_\alpha\) is a group homomorphism,
  \(f_\alpha(u_\alpha)=v_\alpha\), and
  \[
    f_T(t)f_\alpha(x)f_T(t)^{-1}
    =
    f_\alpha(\alpha(t)x)
  \]
  for all \(t\in T(S')\), \(x\in U_\alpha(S')\), and \(S'\to S\);
\item[\SGAInlineRefTarget{sga3:xxiv:lemma:7-1-4:item:ii}(ii)] \(f_\alpha(u_\alpha)=v_\alpha\), and
  \[
    f_\alpha(a(t,t')u_\alpha)=b(t,t')
  \]
  for all \(t,t'\in T(S')\) and \(S'\to S\).
\end{enumerate}}

If \(f_\alpha\) satisfies \SGARefLink{sga3:xxiv:lemma:7-1-4:item:i}{(i)}, then
\[
  f_\alpha(\alpha(t)u_\alpha)
  =
  \operatorname{int}(f_T(t))(v_\alpha),
\]
which immediately gives \SGARefLink{sga3:xxiv:lemma:7-1-4:item:ii}{(ii)}. Conversely, suppose \SGARefLink{sga3:xxiv:lemma:7-1-4:item:ii}{(ii)}. Since \(a\)
is an \textup{(fpqc)} covering, to prove the last condition in \SGARefLink{sga3:xxiv:lemma:7-1-4:item:i}{(i)} it
suffices to show that, for \(t,t',t''\in T(S')\),
\[
  f_T(t)\,
  f_\alpha(a(t',t'')u_\alpha)\,
  f_T(t)^{-1}
  =
  f_\alpha(\alpha(t)a(t',t'')u_\alpha).
\]
This is the same as
\[
  f_T(t)b(t',t'')f_T(t)^{-1}=b(tt',tt''),
\]
which follows directly from the definition of \(b\).

It remains to prove that \(f_\alpha\) is a group homomorphism. The
preceding property gives
\[
\begin{aligned}
 f_\alpha(\alpha(t)u_\alpha)
 f_\alpha(\alpha(t')u_\alpha)
 &=
 \bigl(f_T(t)f_\alpha(u_\alpha)f_T(t)^{-1}\bigr)
 \bigl(f_T(t')f_\alpha(u_\alpha)f_T(t')^{-1}\bigr)\\
 &=b(t,t')\\
  &=f_\alpha\bigl(\alpha(t)u_\alpha+
                   \alpha(t')u_\alpha\bigr).
\end{aligned}
\]
Thus \(f_\alpha(x+x')=f_\alpha(x)f_\alpha(x')\) whenever \(x,x'\) are
sections of the schematically dense open subset
\(U_\alpha^\times\subset U_\alpha\). \SGARefLink{sga3:xviii:prop:1.4}{Exposé XVIII, 1.4} then gives the
result.

\subsubsection*{7.1.5.}

Fix provisionally an \(\alpha\in\Delta\). Since \(a\) is faithfully
flat and quasi-compact, \(G_{a,S}\) is identified with the quotient of
\(T\times_S T\) by the equivalence relation \(R\) defined by \(a\).
Write \(i_1,i_2:R\to T\times_S T\) for its two projections.

\[
  \begin{tikzcd}[column sep=huge]
R
  \arrow[r,shift left=.7ex,"i_1"]
  \arrow[r,shift right=.7ex,"i_2"']
&
T\times_S T
\end{tikzcd}

\]

For \(b\) to factor through \(a\),\footnote{\textit{Editors' note.}
That is, for condition \SGARefLink{sga3:xxiv:lemma:7-1-4:item:ii}{7.1.4(ii)} to hold.} it is necessary and
sufficient that
\[
  b\circ i_1=b\circ i_2.
\]
Let \(K\) be the equalizer of \(b\circ i_1\) and \(b\circ i_2\).

\[
  \begin{tikzcd}[column sep=huge]
R
  \arrow[r,shift left=.7ex,"b\circ i_1"]
  \arrow[r,shift right=.7ex,"b\circ i_2"']
&
H
\end{tikzcd}

\]

Then the required condition is \(K=R\). Since \(H\) is separated over
\(S\), \(K\) is a closed subscheme of \(R\). On the other hand,
\(R\) is essentially free over \(S\), by \SGARefLink{sga3:xxiv:lemma:7-1-3}{7.1.3}.

\subsubsection*{7.1.6.}

It follows that, for an \(S\)-scheme \(S'\), an \(f_\alpha\) satisfying
\SGARefLink{sga3:xxiv:lemma:7-1-4:item:i}{7.1.4(i)} exists over \(S'\), and is then unique, if and only if
\[
  K_{S'}=R_{S'}
\]
and the resulting morphism satisfies
\[
  f_\alpha(u_\alpha)=v_\alpha.
\]
The first condition is equivalent to requiring \(S'\to S\) to factor
through a certain closed subscheme of \(S\) (\SGARefLink{sga3:viii:theorem:6:4}{Exposé VIII,
6.4}\footnote{\textit{Editors' note.} See also the addition in Exposé
VI\textsubscript{B}, \SGARefLink{thm:sga3-VIB-6.2.3}{6.2.3}.}). After replacing \(S\) by this closed
subscheme, the second condition again defines a closed subscheme of
\(S\), since it is the equality of two sections of the separated
\(S\)-scheme \(H\).

After replacing \(S\) by the resulting closed subscheme, we may thus
suppose that a morphism \(f_\alpha:U_\alpha\to H\) satisfying \SGARefLink{sga3:xxiv:lemma:7-1-4:item:i}{7.1.4(i)}
exists. Intersecting the closed subschemes obtained for all
\(\alpha\in\Delta\), we may suppose this condition holds for every
\(\alpha\in\Delta\).

\subsubsection*{7.1.7.}

Similarly, for each \(\alpha\in\Delta\), consider the two
\(S\)-group homomorphisms
\[
  f_T\circ\operatorname{int}(w_\alpha),
  \qquad
  \operatorname{int}(h_\alpha)\circ f_T
  :T\longrightarrow H.
\]
Since \(H\) is separated over \(S\) and \(T\) is essentially free over
\(S\), the same reasoning shows that, after replacing \(S\) by a
closed subscheme, one may suppose that
\[
  f_T\circ\operatorname{int}(w_\alpha)
  =
  \operatorname{int}(h_\alpha)\circ f_T
\]
for every \(\alpha\in\Delta\).

\subsubsection*{7.1.8.}

The generators-and-relations theorem (\SGARefLink{sga3:xxiii:subsection:3-5}{Exposé XXIII, 3.5}) now shows that
a group homomorphism \(f:G\to H\) with the required properties exists
if and only if a certain finite set of algebraic relations among the
sections
\[
  h_\alpha,\qquad
  v_\alpha,\qquad
  f_T\bigl(\alpha^*(-1)\bigr),
  \qquad \alpha\in\Delta,
\]
is satisfied:
\[
  R_i\!\left(
    (h_\alpha)_\alpha,
    (v_\alpha)_\alpha,
    \bigl(f_T(\alpha^*(-1))\bigr)_\alpha
  \right)=e,
  \qquad i=1,\ldots,n.
\]
Because \(H\) is separated over \(S\), these relations again define a
closed subscheme of \(S\), proving \SGARefLink{sga3:xxiv:proposition:7-1-2}{Proposition 7.1.2}.

\par\medskip\noindent\SGAInlineRefTarget{sga3:xxiv:corollary:7-1-9}\textbf{Corollary 7.1.9.---}\enspace
\emph{Let \(S\) be a scheme, \(G\) a split reductive \(S\)-group, \(T\)
its maximal torus, and \(H\) a separated \(S\)-group scheme. Let
\(\mathcal P\) be a property of morphisms satisfying:
\begin{enumerate}
\item[(i)] every closed immersion has \(\mathcal P\);
\item[(ii)] a composite of morphisms having \(\mathcal P\) has
  \(\mathcal P\);
\item[(iii)] \(\mathcal P\) is stable under base change;
\item[(iv)] the structure morphism \(H\to S\) has \(\mathcal P\).
\end{enumerate}
Then the canonical morphism
\[
  \underline{\operatorname{Hom}}_{S\text{-}\mathrm{gr.}}(G,H)
  \longrightarrow
  \underline{\operatorname{Hom}}_{S\text{-}\mathrm{gr.}}(T,H)
\]
is relatively representable by a separated morphism having
\(\mathcal P\).\footnote{\textit{Editors' note.} More explicitly, for
every representable
\(S'\to\underline{\operatorname{Hom}}_{S\text{-}\mathrm{gr.}}(T,H)\),
the fiber product
\(\underline{\operatorname{Hom}}_{S\text{-}\mathrm{gr.}}(G,H)
\times S'\) is representable, and its morphism to \(S'\) is separated
and has \(\mathcal P\).}}

Indeed, we may suppose \(G\) pinned. The structure morphism
\(H^{2n}\to S\) has \(\mathcal P\), and \SGARefLink{sga3:xxiv:proposition:7-1-2}{7.1.2} gives the assertion.

\par\medskip\noindent\SGAInlineRefTarget{sga3:xxiv:corollary:7-1-10}\textbf{Corollary 7.1.10.---}\enspace
\emph{Let \(S\) be a scheme, \(G\) a split reductive \(S\)-group, and
\(H\) a smooth \(S\)-group scheme with affine
fibers.\footnote{\textit{Editors' note.} We have suppressed the
superfluous hypothesis that \(H\) be quasi-separated.} Then
\(\underline{\operatorname{Hom}}_{S\text{-}\mathrm{gr.}}(G,H)\) is
representable by an \(S\)-scheme locally of finite presentation and
separated over \(S\).}

Since \(H\) is smooth over \(S\), its neutral component \(H^0\) is
defined; it has affine fibers and is smooth, separated, and of finite
presentation over \(S\) (Exposé VI\textsubscript{B}, \SGARefLink{thm:sga3-VIB-3.10}{3.10} and \SGARefLink{sga3:xxiv:proposition:5-5}{5.5}).
Since \(G\) has connected fibers, \(H\) may be replaced by \(H^0\).
As \(G\) and \(H\) are then of finite presentation, we may suppose
\(S\) noetherian and apply \SGARefLink{sga3:xxiv:corollary:7-1-9}{7.1.9} with \(\mathcal P\) equal to the
property ``of finite type''. By \SGARefLink{sga3:xv:theorem:8:8}{Exposé XV, 8.9},
\(\underline{\operatorname{Hom}}_{S\text{-}\mathrm{gr.}}(T,H)\) is
representable by a separated \(S\)-scheme locally of finite type.

\par\medskip\noindent\SGAInlineRefTarget{sga3:xxiv:remark:7-1-11}\textbf{Remark 7.1.11.---}\enspace
If \(H\) is affine over \(S\), the reference to Exposé XV may be
replaced by \SGARefLink{sga3:xi:corollary:4:2}{Exposé XI, 4.2}.

\subsection*{\SGAInlineRefTarget{sga3:xxiv:subsection:7-2}7.2. General case}

\par\medskip\noindent\SGAInlineRefTarget{sga3:xxiv:proposition:7-2-1}\textbf{Proposition 7.2.1.---}\enspace
\emph{Let \(S\) be a scheme, \(G\) a reductive \(S\)-group, and \(T\) a
maximal torus of \(G\). Let \(H\) be an \(S\)-group scheme whose
structure morphism \(H\to S\) satisfies:
\begin{quote}
\textup{(+)} Every point \(s\in S\) has an open neighborhood \(U\)
such that \(H_U\to U\) is quasi-projective.
\end{quote}
Then the canonical morphism
\[
  \underline{\operatorname{Hom}}_{S\text{-}\mathrm{gr.}}(G,H)
  \longrightarrow
  \underline{\operatorname{Hom}}_{S\text{-}\mathrm{gr.}}(T,H)
\]
is relatively representable by a morphism satisfying \textup{(+)}.}

When \(G\) is split relative to \(T\), apply 7.1.9 with
\(\mathcal P=\textup{(+)}\). When \(G\) is locally isotrivial, for
example when \(G\) is semisimple (4.2.2), note that the assertion is
local for the local finite étale topology. Indeed,
quasi-projectivity is local for the global finite étale topology and
ensures effectiveness of descent for that topology (SGA~1, VIII, 7.7).
The general case follows from the next lemma.

\par\medskip\noindent\SGAInlineRefTarget{sga3:xxiv:lemma:7-2-2}\textbf{Lemma 7.2.2.---}\enspace
\emph{Let \(S\) be a scheme, \(G\) a reductive \(S\)-group, \(T\) a
maximal torus of \(G\), \(G'\) the derived group of \(G\) (Exposé XXII,
6.2), and \(T'=T\cap G'\) the maximal torus of \(G'\) corresponding to
\(T\) (Exposé XXII, 6.2.8). The natural diagram involving
\(G,G',T,T'\) is a pushout in the category of \(S\)-group sheaves.
Equivalently, for every \(S\)-group sheaf \(H\), the corresponding
square of Hom-functors is cartesian.}

\[
  \begin{tikzcd}[column sep=large,row sep=large]
G & G' \arrow[l] \\
T \arrow[u] & T' \arrow[l] \arrow[u]
\end{tikzcd}

\]

\[
  \begin{tikzcd}[column sep=large,row sep=large]
\underline{\operatorname{Hom}}_{S\text{-}\mathrm{gr.}}(G,H)
  \arrow[r]
  \arrow[d]
&
\underline{\operatorname{Hom}}_{S\text{-}\mathrm{gr.}}(G',H)
  \arrow[d]
\\
\underline{\operatorname{Hom}}_{S\text{-}\mathrm{gr.}}(T,H)
  \arrow[r]
&
\underline{\operatorname{Hom}}_{S\text{-}\mathrm{gr.}}(T',H)
\end{tikzcd}

\]

If \(\operatorname{rad}(G)\) is the radical of \(G\), then
\(\operatorname{rad}(G)\subset T\), so
\[
  \operatorname{rad}(G)\cap G'
  =
  \operatorname{rad}(G)\cap T'
  =
  K.
\]
Multiplication in \(G\) induces isogenies (\SGARefLink{sga3:xxii:subsection:6-2}{Exposé XXII, 6.2})
\[
  i:G'\times_S\operatorname{rad}(G)\longrightarrow G,
  \qquad
  j:T'\times_S\operatorname{rad}(G)\longrightarrow T.
\]
Let
\[
  f_{G'}:G'\to H,\qquad
  f_T:T\to H,\qquad
  f_{T'}:T'\to H
\]
be \(S\)-group homomorphisms such that
\[
  f_{G'}|_{T'}=f_T|_{T'}=f_{T'}.
\]
We show that there is a unique \(S\)-group homomorphism \(f:G\to H\)
inducing \(f_{G'}\) and \(f_T\). Put
\[
  f_{\mathrm{rad}}
  =
  f_T|_{\operatorname{rad}(G)}
\]
\footnote{\textit{Editors' note.} We have replaced \(f_r\) by
\(f_{\mathrm{rad}}\).}
and define
\[
  f_1=f_{G'}\cdot f_{\mathrm{rad}}:
  G'\times_S\operatorname{rad}(G)\longrightarrow H.
\]
For \(f\) to exist, and it is plainly then unique, it is necessary and
sufficient that \(f_1\) be trivial on the kernel of \(i\), that is,
that \(f_{G'}\) and \(f_{\mathrm{rad}}\) coincide on \(K\). The
existence of \(f_T\) similarly shows that \(f_{T'}\) and
\(f_{\mathrm{rad}}\) coincide on \(K\). Finally, \(f_{G'}\) and
\(f_{T'}\) coincide on \(K\subset T'\), proving the lemma.

Arguing as in \SGARefLink{sga3:xxiv:corollary:7-1-10}{7.1.10}, \SGARefLink{sga3:xxiv:proposition:7-2-1}{Proposition 7.2.1} gives:

\par\medskip\noindent\SGAInlineRefTarget{sga3:xxiv:corollary:7-2-3}\textbf{Corollary 7.2.3.---}\enspace
\emph{Let \(S\) be a scheme, \(G\) a reductive \(S\)-group, and \(H\)
a smooth, quasi-projective \(S\)-group scheme with affine fibers. Then
\(\underline{\operatorname{Hom}}_{S\text{-}\mathrm{gr.}}(G,H)\) is
representable by an \(S\)-scheme locally of finite presentation and
separated over \(S\).}

\subsection*{\SGAInlineRefTarget{sga3:xxiv:subsection:7-3}7.3. Phenomena specific to characteristic zero}

If \(G\) and \(H\) are smooth \(S\)-groups, with Lie algebras
\(\mathfrak g\) and \(\mathfrak h\), there is a canonical morphism
\[
  \Lie:
  \underline{\operatorname{Hom}}_{S\text{-}\mathrm{gr.}}(G,H)
  \longrightarrow
  \underline{\operatorname{Hom}}_{\mathcal O_S\text{-}\mathrm{Lie}}
    (\mathfrak g,\mathfrak h),
\]
where the functor on the right has the evident meaning.

\par\medskip\noindent\SGAInlineRefTarget{sga3:xxiv:proposition:7-3-1}\textbf{Proposition 7.3.1.---}\enspace
\emph{Let \(S\) be a scheme of characteristic zero, that is, a
\(\mathbb Q\)-scheme; let \(G\) be a reductive \(S\)-group, and \(H\) a
smooth, quasi-projective \(S\)-group scheme with affine fibers.
\begin{enumerate}
\item[\SGAInlineRefTarget{sga3:xxiv:proposition:7-3-1:item:i}(i)]
\(\underline{\operatorname{Hom}}_{S\text{-}\mathrm{gr.}}(G,H)\) is
representable by a smooth separated \(S\)-scheme.
\item[(ii)] If \(G\) is semisimple, this scheme is affine and of finite
presentation over \(S\).
\item[\SGAInlineRefTarget{sga3:xxiv:proposition:7-3-1:item:iii}(iii)] If \(G\) is simply connected (\SGARefLink{sga3:xxii:def:4.3.3}{Exposé XXII, 4.3.3}), the
canonical morphism
\[
  \Lie:
  \underline{\operatorname{Hom}}_{S\text{-}\mathrm{gr.}}(G,H)
  \longrightarrow
  \underline{\operatorname{Hom}}_{\mathcal O_S\text{-}\mathrm{Lie}}
    (\mathfrak g,\mathfrak h)
\]
is an isomorphism.
\end{enumerate}}

\par\medskip\noindent\SGAInlineRefTarget{sga3:xxiv:lemma:7-3-2}\textbf{Lemma 7.3.2.---}\enspace
\emph{Let \(k\) be a field of characteristic zero, \(G\) a reductive
\(k\)-group, \(V\) a finite-dimensional \(k\)-vector space, and
\(G\to\operatorname{GL}(V)\) a linear representation. Then
\[
  H^0(G,V)=H^0(\mathfrak g,V),
  \qquad
  H^i(G,V)=0\quad(i>0).
\]}

The first equality is valid in general for a smooth connected
group.\footnote{\textit{Editors' note.} The inclusion
\(V^G\subset V^{\mathfrak g}\) is clear. Conversely,
\(C=\operatorname{Centr}_G(V^{\mathfrak g})\) is a closed subgroup of
\(G\), smooth because \(\operatorname{char}(k)=0\). By \SGARefLink{sga3:ii:theorem:5-3-1}{Exposé II,
5.3.1},
\(\Lie(C)=\operatorname{Centr}_{\mathfrak g}(V^{\mathfrak g})
=\mathfrak g\); because \(C\) is smooth and \(G\) connected, \(C=G\).
Thus \(V^{\mathfrak g}\subset V^G\). See also \([{\rm DG70}]\),
\S II.6, Proposition 2.1(c).} For a reductive group it can also be
proved as follows. We may suppose \(k\) algebraically closed, hence
\(G\) split and generated by subgroups isomorphic to
\(G_{m,k}\),\footnote{\textit{Editors' note.} The union of the maximal
tori of \(G\) is dense in \(G\); cf.\ \emph{Bible}, \S6.5,
Theorem~5, or \([{\rm Ch05}]\), \S6.6, Theorem~6.} and the assertion
for this group is immediate.

The first equality shows that the functor \(H^0(G,V)\) is exact in
\(V\) when \(G\) is semisimple: \(\mathfrak g\) is then a semisimple
Lie algebra, and one applies \([{\rm BLie}]\), \S I.6,
Exercise~1(b). The assertion remains true when \(G\) is reductive.
Indeed, introducing the radical \(R\) of
\(G\)\footnote{\textit{Editors' note.} Observe that \(R\) is a torus.}
and the semisimple quotient \(G/R\), one has
\[
  H^0(G,V)=H^0\bigl(G/R,H^0(R,V)\bigr),
\]
so \(H^0(G,-)\) is a composite of two exact functors. \SGARefLink{sga3:i:theorem:5-3-1}{Exposé I, 5.3.1}
then gives \(H^i(G,V)=0\) for \(i>0\).

\par\medskip\noindent\SGAInlineRefTarget{sga3:xxiv:remark:7-3-3}\textbf{Remark 7.3.3.---}\enspace
Under the preceding hypotheses, if \(G\) is semisimple, then
\[
  H^1(\mathfrak g,V)=H^2(\mathfrak g,V)=0;
\]
see \([{\rm BLie}]\), \emph{loc.\ cit.}, (b) and (d).

\subsubsection*{7.3.4.}

We prove the proposition. By \SGARefLink{sga3:xxiv:corollary:7-2-3}{7.2.3},
\(\underline{\operatorname{Hom}}_{S\text{-}\mathrm{gr.}}(G,H)\) is
already representable by an \(S\)-scheme locally of finite
presentation and separated over \(S\). To show it smooth, it suffices
to show it infinitesimally smooth (\SGARefLink{sga3:xi:remark:1:8}{Exposé XI,
1.8}),\footnote{\textit{Editors' note.} This means that for every local
artinian \(S\)-scheme \(S'\), with closed point \(s'\), every
homomorphism of \(\kappa(s')\)-groups
\(G_{s'}\to H_{s'}\) lifts to an \(S'\)-group homomorphism
\(G_{S'}\to H_{S'}\). By \SGARefLink{sga3:iii:corollary:2-8}{Exposé III, 2.8}, this follows from the
vanishing of \(H^2(G_{s'},V)\), where
\(V=\Lie(H_{s'}/\kappa(s'))\).} and this follows from \SGARefLink{sga3:iii:corollary:2-8:item:i}{Exposé III,
2.8(i)} and \SGARefLink{sga3:xxiv:lemma:7-3-2}{7.3.2}. This proves \SGARefLink{sga3:xxiv:proposition:7-3-1:item:i}{(i)}.

We next show that \SGARefLink{sga3:iii:corollary:2-8:item:ii}{(ii)} follows from \SGARefLink{sga3:xxiv:proposition:7-3-1:item:iii}{(iii)}. First it is enough to prove
that
\(\underline{\operatorname{Hom}}_{S\text{-}\mathrm{gr.}}(G,H)\) is
affine over \(S\), since it is then of finite presentation, being
smooth by \SGARefLink{sga3:xxiv:proposition:7-3-1:item:i}{(i)}. In any case,
\[
  \underline{\operatorname{Hom}}_{\mathcal O_S\text{-}\mathrm{Lie}}
    (\mathfrak g,\mathfrak h)
\]
is represented by a closed subscheme of
\[
  \underline{\operatorname{Hom}}_{\mathcal O_S\text{-}\mathrm{mod}}
    (\mathfrak g,\mathfrak h)
  \simeq
  \mathrm W\!\left(
    \mathfrak g^\vee\otimes_{\mathcal O_S}\mathfrak h
  \right),
\]
which is an affine \(S\)-scheme. Thus the desired conclusion follows
when \(G\) is simply connected.

In general, we may suppose \(G\) split, so that
\(G\simeq\operatorname{\acute Ep}_S(\mathscr R)\). Introduce the
simply connected root datum \(\operatorname{sc}(\mathscr R)\)
(\SGARefLink{sga3:XXI:6.5.5}{Exposé XXI, 6.5.5}). The existence theorem (\SGARefLink{sga3:xxv:theorem:1:1}{Exposé XXV, 1.1}) gives a
simply connected \(S\)-group \(\widetilde G\) and a central isogeny
\[
  \widetilde G\longrightarrow G.
\]
Its kernel \(K\) is a finite diagonalizable \(S\)-group, hence a
twisted constant \(S\)-group because \(S\) has characteristic zero.
The functor
\(\underline{\operatorname{Hom}}_{S\text{-}\mathrm{gr.}}(K,H)\) is
plainly represented by a separated \(S\)-scheme. For example, if
\(K\simeq(\mathbb Z/n\mathbb Z)_S\), then it is isomorphic to
\[
  \ker\!\left(H\xrightarrow{\,n\,}H\right).
\]
There is an exact sequence of pointed \(S\)-schemes
\[
  1\longrightarrow
  \underline{\operatorname{Hom}}_{S\text{-}\mathrm{gr.}}(G,H)
  \xrightarrow{u}
  \underline{\operatorname{Hom}}_{S\text{-}\mathrm{gr.}}
    (\widetilde G,H)
  \longrightarrow
  \underline{\operatorname{Hom}}_{S\text{-}\mathrm{gr.}}(K,H).
\]
Thus \(u\) is a closed immersion, and
\(\underline{\operatorname{Hom}}_{S\text{-}\mathrm{gr.}}(G,H)\) is
affine over \(S\).

\SGAInlineRefTarget{sga3:xxiv:subsubsection:7-3-5}\subsubsection*{7.3.5.}

It remains to prove \SGARefLink{sga3:xxiv:proposition:7-3-1:item:iii}{(iii)}. As in the proof of \SGARefLink{sga3:xxiv:proposition:7-3-1:item:i}{(i)}, and using \SGARefLink{sga3:xxiv:remark:7-3-3}{7.3.3},
one shows that
\[
  \underline{\operatorname{Hom}}_{\mathcal O_S\text{-}\mathrm{Lie}}
    (\mathfrak g,\mathfrak h)
\]
is smooth over \(S\). We may therefore suppose
\(S=\operatorname{Spec}(k)\), where \(k\) is algebraically closed of
characteristic zero. It is enough to show that \(\Lie\) is bijective
on \(k\)-points and induces a bijection on tangent spaces at
corresponding points. First we show that
\[
  \operatorname{Hom}_{k\text{-}\mathrm{gr.}}(G,H)
  \longrightarrow
  \operatorname{Hom}_{k\text{-}\mathrm{Lie}}
    (\mathfrak g,\mathfrak h)
\]
is bijective.

If \(u:G\to H\) is a \(k\)-group homomorphism, its graph is a connected
subgroup of \(G\times_k H\), determined by its Lie algebra, which is
the graph of \(\Lie(u)\). Hence the map is injective. Conversely, if
\(v:\mathfrak g\to\mathfrak h\) is a Lie algebra homomorphism, its
graph \(\mathfrak k\) is a Lie subalgebra of
\(\mathfrak g\times\mathfrak h\), isomorphic to \(\mathfrak g\).
Because \(\mathfrak g\) is its own derived algebra, so is
\(\mathfrak k\). A theorem of Chevalley
\([{\rm Ch51}]\), \S14, Theorem~15, therefore shows that
\(\mathfrak k\) is the Lie algebra of a connected subgroup \(K\) of
\(G\times_k H\).\footnote{\textit{Editors' note.} See also
\([{\rm Bo91}]\), II 7.9. We have also added the following sentence.}

Since \(\mathfrak k\simeq\mathfrak g\) is semisimple, \(K\) is a
semisimple \(k\)-group (the French prints \(G\) here; the argument
and the following sentence require \(K\)).
Moreover,
\[
  \dim K=\dim\mathfrak k=\dim\mathfrak g=\dim G,
\]
and \(K\cap H\) is finite because its Lie algebra is zero. Thus the
canonical morphism
\[
  \operatorname{pr}_1:K\longrightarrow G
\]
is finite and dominant; because \(G\) is connected, it is surjective,
and hence an isogeny. Since \(G\) is simply connected and \(K\) is
semisimple, \SGARefLink{sga3:XXI:6.2.7}{Exposé XXI, 6.2.7} shows that \(\operatorname{pr}_1\) is an
isomorphism. Consequently \(K\) is the graph of a \(k\)-group
homomorphism \(u:G\to H\) with \(\Lie(u)=v\).

Finally, let \(u:G\to H\) be a \(k\)-group homomorphism. The tangent
space at \(u\) to
\[
  \underline{\operatorname{Hom}}_{k\text{-}\mathrm{gr.}}(G,H)
\]
identifies with \(Z^1(G,\mathfrak h)\) (\SGARefLink{sga3:ii:subsection:4-2}{Exposé II, 4.2}), where
\(G\) acts on \(\mathfrak h\) through
\(\operatorname{Ad}_H\circ u\). Similarly, the tangent space to
\[
  \underline{\operatorname{Hom}}_{k\text{-}\mathrm{Lie}}
    (\mathfrak g,\mathfrak h)
\]
at \(\Lie(u)\) identifies with \(Z^1(\mathfrak g,\mathfrak h)\). We
must show that
\[
  Z^1(G,\mathfrak h)\longrightarrow
  Z^1(\mathfrak g,\mathfrak h)
\]
is bijective.

\[
  \begin{tikzcd}[column sep=small,row sep=large]
0 \arrow[r]
&
\mathfrak h^G
  \arrow[r]
  \arrow[d]
&
\mathfrak h
  \arrow[r]
  \arrow[d,equal]
&
Z^1(G,\mathfrak h)
  \arrow[r]
  \arrow[d]
&
H^1(G,\mathfrak h)
  \arrow[d]
\\
0 \arrow[r]
&
\mathfrak h^{\mathfrak g}
  \arrow[r]
&
\mathfrak h
  \arrow[r]
&
Z^1(\mathfrak g,\mathfrak h)
  \arrow[r]
&
H^1(\mathfrak g,\mathfrak h)
\end{tikzcd}

\]

By \SGARefLink{sga3:xxiv:lemma:7-3-2}{7.3.2} and \SGARefLink{sga3:xxiv:remark:7-3-3}{7.3.3},
\[
  \mathfrak h^G=\mathfrak h^{\mathfrak g},
  \qquad
  H^1(G,\mathfrak h)=0=H^1(\mathfrak g,\mathfrak h),
\]
which gives the desired conclusion.

\par\medskip\noindent\SGAInlineRefTarget{sga3:xxiv:remark:7-3-6}\textbf{Remark 7.3.6.---}\enspace
It is presumably true that, if \(S\) is locally noetherian, \(G\) a
simply connected semisimple \(S\)-group, \(H\) a smooth \(S\)-group
scheme, and \(\widehat G,\widehat H\) their formal completions along
the identity section, then every homomorphism
\(\widehat G\to\widehat H\) comes from a unique homomorphism \(G\to H\).
This would generalize \SGARefLink{sga3:xxiv:proposition:7-3-1:item:iii}{7.3.1(iii)}. When \(S=\operatorname{Spec}(k)\)
and \(H\) is affine and of finite type, this follows from a theorem of
Dieudonné \([{\rm Di57}]\), \S15, Theorem~4.\footnote{\textit{Editors'
note.} See Exposé VII\textsubscript{B} for the definition of the formal
groups \(\widehat G\) and \(\widehat H\). Suppose that
\(S=\operatorname{Spec}(k)\), where \(k\) is a field. By loc. cit.,
2.2.1, giving a morphism of formal \(k\)-groups
\(v:\widehat G\to\widehat H\) is equivalent to giving a morphism of
Hopf \(k\)-algebras
\(\phi:\operatorname{Dist}(G)\to\operatorname{Dist}(H)\), where
\(\operatorname{Dist}(G)\) denotes the algebra of distributions (at
the origin) of \(G\); see Exposé VII\textsubscript{A}, \SGARefLink{sga3:VIIA:2.1}{2.1},
\([{\rm DG70}]\), \S II.4, \SGARefLink{sga3:xxiv:lemma:6-1}{6.1}, or \([{\rm Ja87}]\), I 7.7. It is called
the hyperalgebra of \(G\) in \([{\rm Di57}]\) and \([{\rm Ta75}]\).
Theorem~4 of \([{\rm Di57}]\) (see also \([{\rm Ta75}]\), 0.3.4 (f)
and (g)) generalizes in this context the theorem of Chevalley used in
\SGARefLink{sga3:xxiv:subsubsection:7-3-5}{7.3.5}: one obtains a closed connected \(k\)-subgroup
\(K\subset G\times H\) such that \(\widehat K\) equals the graph of
\(v\). Since \(\operatorname{Dist}(K)\to\operatorname{Dist}(G)\) is an
isomorphism, \(K\to G\) is finite étale. It follows that \(K\) is
semisimple and then, by \SGARefLink{sga3:XXI:6.2.7}{Exposé XXI, 6.2.7}, that \(K\to G\) is an
isomorphism (because \(G\) is simply connected; see also
\([{\rm Ta75}]\), \SGARefLink{sga3:xxiv:1.8}{1.8} and \SGARefLink{sga3:xxiv:2.2}{2.2}). More generally, loc. cit. studies the
\(k\)-groups \(G\) having property \textup{(SC)}: every finite étale
morphism of \(k\)-groups \(G'\to G\), with \(G'\) connected, is an
isomorphism. Finally, the preceding argument shows that every
finite-dimensional \(\operatorname{Dist}(G)\)-module is, uniquely, a
\(G\)-module; for an extension to a split reductive \(k\)-group \(G\)
(or even to a Borel subgroup of \(G\)), see \([{\rm Ja87}]\), II 1.20
(and the references given there).}

\subsection*{\SGAInlineRefTarget{sga3:xxiv:subsection:7-4}7.4. An example}

As an example, we shall determine
\[
  \underline{\operatorname{Hom}}_{S\text{-}\mathrm{gr.}}
    (\mathrm{SL}_{2,S},\mathrm{SL}_{2,S}).
\]

\subsubsection*{7.4.1.}

Recall (Exposé XX, no.~5) that \(\mathrm{SL}_{2,S}\) is the
\(S\)-group scheme of matrices
\[
  \begin{pmatrix}a&b\\c&d\end{pmatrix}
\]
over \(S\) satisfying \(ad-bc=1\). A maximal torus \(T\) is defined
by the monomorphism
\[
  \alpha^*:\mathbb G_{m,S}\longrightarrow\mathrm{SL}_{2,S},
  \qquad
  \alpha^*(z)=
  \begin{pmatrix}z&0\\0&z^{-1}\end{pmatrix}.
\]
A root of \(G\) with respect to \(T\) is defined by
\[
  \alpha\bigl(\alpha^*(z)\bigr)=z^2,
\]
and a corresponding monomorphism is
\[
  p:\mathbb G_{a,S}\longrightarrow\mathrm{SL}_{2,S},
  \qquad
  p(x)=
  \begin{pmatrix}1&x\\0&1\end{pmatrix}.
\]
Finally, putting \(u=p(1)\), the corresponding representative of the
Weyl group is
\[
  w=
  \begin{pmatrix}0&1\\-1&0\end{pmatrix}.
\]

Recall also (\SGARefLink{sga3:xx:6.2}{Exposé XX, 6.2}) that \(\mathrm{SL}_{2,S}\) is generated by
\(T\), \(p(\mathbb G_{a,S})\), and \(w\), subject to the relations
\[
  \alpha^*(z)p(x)\alpha^*(z^{-1})=p(z^2x),
\]
\[
  w\alpha^*(z)w^{-1}=\alpha^*(z^{-1}),
  \qquad
  w^2=\alpha^*(-1),
  \qquad
  (wu)^3=1.
\]

\subsubsection*{7.4.2.}

Let \(f\) be an endomorphism of \(\mathrm{SL}_{2,S}\). It first
defines a homomorphism
\[
  f\circ\alpha^*:\mathbb G_{m,S}\longrightarrow\mathrm{SL}_{2,S}.
\]
The closed subgroup
\(\ker(f\circ\alpha^*)\subset\mathbb G_{m,S}\) is, locally on \(S\),
either \(\mu_{n,S}\), with \(n\geq1\), or
\(\mathbb G_{m,S}\). After restricting \(S\), we may therefore assume
that there are an integer \(n\in\mathbb N\) and a monomorphism
\[
  f':\mathbb G_{m,S}\longrightarrow\mathrm{SL}_{2,S}
\]
such that
\[
  f\circ\alpha^*(z)=f'(z^n).
\]
By conjugacy of monomorphisms
\(\mathbb G_{m,S}\to\mathrm{SL}_{2,S}\), after an étale covering
extension of the base there is a section \(g\) of \(G\) such that
\[
  f\circ\alpha^*(z)
  =
  \operatorname{int}(g)\circ\alpha^*(z^n).
\]
Replacing \(f\) by an inner conjugate, we are thus reduced to the case
\[
  f\circ\alpha^*(z)=\alpha^*(z^n)
\]
for some \(n\in\mathbb N\).

\subsubsection*{7.4.3.}

Consider now
\[
  f\circ p:\mathbb G_{a,S}\longrightarrow\mathrm{SL}_{2,S}.
\]
It satisfies
\[
  \alpha^*(z^n)(f\circ p)(x)\alpha^*(z^n)^{-1}
  =
  (f\circ p)(z^2x).
\]
If \(n=0\), it follows at once that \(f\circ p\) is invariant under
the homotheties of \(\mathbb G_{a,S}\), and hence is constant. If
\(n\neq0\), \SGARefLink{sga3:xxii:prop:4.1.9}{Exposé XXII, 4.1.9} applies: the map
\(x\mapsto x^n\) is an endomorphism of \(\mathbb G_{a,S}\), and there
is a \(\lambda\in\mathbb G_{m,S}\) such that
\[
  f\circ p(x)=p(\lambda x^n).
\]
This relation is also valid for \(n=0\), on taking \(\lambda=1\).
After extending \(S\) once more, choose
\(z\in\mathbb G_{m,S}\) with \(z^2=\lambda\). Replacing \(f\) by
\(\operatorname{int}(\alpha^*(z))\circ f\), we are reduced to
\[
  f\circ\alpha^*(z)=\alpha^*(z^n),
  \qquad
  f\circ p(x)=p(x^n).
\]

\subsubsection*{7.4.4.}

Finally, one must have
\[
  f(w)\alpha^*(z^n)f(w)^{-1}=\alpha^*(z^n)^{-1},
  \qquad
  (f(w)u)^3=1.
\]
By \SGARefLink{sga3:xx:3.8}{Exposé XX, 3.8}, this implies \(f(w)=w^n\). Since \(G\) is generated
by \(T\), \(p(\mathbb G_{a,S})\), and \(w\), it follows that, for every
\(S'\to S\) and every
\(\begin{pmatrix}a&b\\c&d\end{pmatrix}\in G(S')\),
\[
  f\!\begin{pmatrix}a&b\\c&d\end{pmatrix}
  =
  \begin{pmatrix}a^n&b^n\\c^n&d^n\end{pmatrix}.
\]

\subsubsection*{7.4.5.}

Summarizing, locally on \(S\) for the étale topology, every
endomorphism \(f\) of \(\mathrm{SL}_{2,S}\) can be written
\[
  f=\phi\circ F_n,
\]
where \(\phi\) is an inner automorphism, \(n\geq0\), and
\[
  F_n\!\begin{pmatrix}a&b\\c&d\end{pmatrix}
  =
  \begin{pmatrix}a^n&b^n\\c^n&d^n\end{pmatrix}.
\]
If \(f=\phi\circ F_n\), the integer \(n\) is determined by any fiber
\(f_s\), for example by \(\ker(f_s)\). Consequently \(n\) is a
locally constant function on \(S\).

\subsubsection*{7.4.6.}

It follows that, for every endomorphism \(f\) of
\(\mathrm{SL}_{2,S}\), the scheme \(S\) has a canonical decomposition
into open and closed subschemes
\[
  S_0,\qquad S_1,\qquad S_{p^n},
\]
where \(p^n\) runs through the positive powers of prime numbers, such
that:
\begin{enumerate}
\item[(i)] \(f_{S_0}\) is the zero morphism;
\item[(ii)] \(f_{S_1}\) is an isomorphism, necessarily an inner
  automorphism;
\item[(iii)] \(S_{p^n}\) has characteristic \(p\), and
  \(f_{S_{p^n}}\) has a unique decomposition
  \[
    f_{S_{p^n}}=\phi\circ F_p^n,
  \]
  where \(\phi\) is an inner automorphism and \(F_p\) is the Frobenius
  endomorphism of \(\mathrm{SL}_{2,\mathbb F_p}\).
\end{enumerate}

\subsubsection*{7.4.7.}

Equivalently,
\[
  \underline{\operatorname{Hom}}_{\mathbb Z\text{-}\mathrm{gr.}}
    (\mathrm{SL}_{2,\mathbb Z},\mathrm{SL}_{2,\mathbb Z})
\]
is the disjoint-union scheme consisting of:
\begin{enumerate}
\item[(i)] one scheme isomorphic to \(\operatorname{Spec}(\mathbb Z)\);
\item[(ii)] one scheme isomorphic to
  \[
    \underline{\operatorname{Aut}}_{\mathbb Z\text{-}\mathrm{gr.}}
      (\mathrm{SL}_{2,\mathbb Z})
    \simeq
    \operatorname{ad}(\mathrm{SL}_{2,\mathbb Z});
  \]
\item[(iii)] for each prime \(p\) and each integer \(n>0\), one scheme
  isomorphic to
  \[
    \underline{\operatorname{Aut}}_{\mathbb F_p\text{-}\mathrm{gr.}}
      (\mathrm{SL}_{2,\mathbb F_p})
    \simeq
    \operatorname{ad}(\mathrm{SL}_{2,\mathbb F_p}).
  \]
\end{enumerate}

\subsubsection*{7.4.8.}

In particular:
\begin{enumerate}
\item[(i)]
\(\underline{\operatorname{Hom}}_{\mathbb F_p\text{-}\mathrm{gr.}}
(\mathrm{SL}_{2,\mathbb F_p},\mathrm{SL}_{2,\mathbb F_p})\)
has infinitely many connected components and is therefore not
quasi-compact.
\item[(ii)] If \(S\) is a scheme of unequal characteristics, then
\(\underline{\operatorname{Hom}}_{S\text{-}\mathrm{gr.}}
(\mathrm{SL}_{2,S},\mathrm{SL}_{2,S})\) is not flat over \(S\).
\end{enumerate}

\section*{\SGAInlineRefTarget{sga3:xxiv:section:8}8. Appendix: Cohomology of a smooth group\\
over a henselian ring. Cohomology and the functor \(\prod\)}

\par\medskip\noindent\SGAInlineRefTarget{sga3:xxiv:proposition:8-1}\textbf{Proposition 8.1.---}\enspace
\emph{Let \(S\) be a henselian local scheme, \(s\) its closed point,
and \(G\) a smooth \(S\)-group scheme such that every finite subset of
\(G\) is contained in an affine open subset.%
\begingroup
\renewcommand{\thefootnote}{*}%
\footnote{This last condition is in fact unnecessary; see
A.~Grothendieck, \textit{Le groupe de Brauer III}, in
\textit{Dix Exposés sur la cohomologie des schémas}, North-Holland,
1968, Theorem~11.7 and Remarks~11.8.3.}%
\endgroup
\addtocounter{footnote}{-1}
\begin{enumerate}
\item[(i)] If \(P\) is a principal homogeneous bundle under \(G\),
  there is a finite surjective étale morphism \(S'\to S\) which
  trivializes \(P\). Moreover,
  \[
    \operatorname{Fib}(S,G)\simeq H^1_{\acute et}(S,G).
  \]
\item[\SGAInlineRefTarget{sga3:xxiv:proposition:8-1:item:ii}(ii)] For every finite surjective étale morphism \(S'\to S\),
  the canonical map
  \[
    H^1(S'/S,G)
    \longrightarrow
    H^1(S'\otimes_S\kappa(s)/\kappa(s),G_s)
  \]
  is bijective.
\item[(iii)] The canonical map
  \[
    \operatorname{Fib}(S,G)
    \longrightarrow
    \operatorname{Fib}(\kappa(s),G_s)
  \]
  is bijective.
\end{enumerate}}

\subsubsection*{8.1.2.}

If \(K\) is a finite separable extension of \(\kappa(s)\), there is a
finite surjective étale morphism \(S'\to S\) such that
\[
  K\simeq S'\otimes_S\kappa(s).
\]
\footnote{\textit{Editors' note.} The scheme \(S'\) is still local and
henselian. We have also retained the numbering of the original: there
is no no.~8.1.1.}
If \(P\) is a principal homogeneous bundle under \(G\), then \(P\) is
smooth over \(S\), hence \(P_s\) is smooth over \(\kappa(s)\).
Consequently there is a finite separable extension \(K/\kappa(s)\)
such that \(P_K\) has a section (EGA IV\textsubscript{4}, 17.15.10).
Representing \(K\) as above, Hensel's lemma
(EGA IV\textsubscript{4}, 18.5.17) shows that \(P_{S'}\) has a section.
This proves the first assertion of (i).

Conversely, if \(P\) is a principal homogeneous sheaf under \(G\) for
the étale topology, then a finite surjective étale morphism
\(S'\to S\) trivializes \(P\): every étale covering family of a
henselian local scheme is dominated by a family consisting of one
morphism of that form. By the descent hypothesis imposed on \(G\),
\(P\) is representable (SGA~1, VIII, 7.6). This completes the proof of
(i), and shows that (ii) implies (iii). It remains to prove (ii).

\subsubsection*{8.1.3.}

The map in (ii) is injective. Let \(P\) and \(Q\) be two principal
homogeneous bundles under \(G\), trivialized by \(S'\), and consider
the \(S\)-group sheaf
\[
  \mathcal H=
  \underline{\operatorname{Isom}}_{G\text{-}\mathrm{bundles}}(P,Q).
\]
Since \(\mathcal H_{S'}\) is isomorphic to \(G_{S'}\),
\(\mathcal H\) is representable by the second hypothesis on \(G\), as
above. If \(\mathcal H(\kappa(s))\neq\varnothing\), Hensel's lemma
gives \(\mathcal H(S)\neq\varnothing\), and hence \(P\simeq Q\).

\subsubsection*{8.1.4.}

To prove surjectivity in \SGARefLink{sga3:xxiv:proposition:8-1:item:ii}{(ii)}, it is equivalent to prove that
\[
  Z^1(S'/S,G)
  \longrightarrow
  Z^1(S'\otimes_S\kappa(s)/\kappa(s),G_s)
\]
is surjective. Let \(Z^1\) be the \(S\)-functor defined by
\[
  Z^1(T)=Z^1(S'\times_ST/T,G_T).
\]
The preceding map is then
\[
  Z^1(S)\longrightarrow Z^1(\kappa(s)).
\]
Another application of Hensel's lemma reduces the assertion to
showing that \(Z^1\) is representable by a smooth \(S\)-scheme.

\subsubsection*{8.1.5.}

We first show that \(Z^1\) is representable by an \(S\)-scheme locally
of finite presentation. For \(i=0,1,\ldots\), define the \(S\)-functor
\(C_i\) by
\[
  C_i(T)=C_i(S'\times_ST/T,G_T).
\]
Equivalently,
\[
  C_i(T)
  =
  G\bigl((S'/S)^{i+1}\times_ST\bigr),
  \qquad
  C_i
  =
  \underline{\operatorname{Hom}}_S((S'/S)^{i+1},G).
\]
Since \(Z^1\) is obtained from \(C_1\) and \(C_2\) by fiber products,
it is enough to prove that \(C_i\), for \(i=1,2\), is representable by
an \(S\)-scheme locally of finite presentation.

\subsubsection*{8.1.6.}

Let \(T\to S\) be a finite surjective étale morphism which splits
\(S'\). Then
\[
  C_i\times_ST
  =
  \underline{\operatorname{Hom}}_T
    ((S'\times_ST/T)^{i+1},G_T)
\]
is represented by a product of \(n\) copies of \(G_T\), where \(n\)
is the degree of \((S'/S)^{i+1}\). The hypothesis on \(G\), applied
once again, shows that \(C_i\) is represented by an \(S\)-scheme
locally of finite presentation (SGA~1, VIII, \textit{loc. cit.}).

\subsubsection*{8.1.7.}

It remains to prove that \(Z^1\) is smooth. By definition, if \(T\) is
an affine \(S\)-scheme and \(T_0\) is the closed subscheme defined by
a square-zero ideal \(J\), we must show that
\[
  Z^1(T)\longrightarrow Z^1(T_0)
\]
is surjective. Since \(G\) is smooth, the canonical map
\(G(T)\to G(T_0)\) is surjective, and it is enough to prove that
\[
  H^1(S'\times_ST/T,G_T)
  \longrightarrow
  H^1(S'\times_ST_0/T_0,G_{T_0})
\]
is bijective. After a slight change of notation and a harmless
generalization of the hypotheses, this reduces to the following
lemma.

\par\medskip\noindent\SGAInlineRefTarget{sga3:xxiv:lemma:8-1-8}\textbf{Lemma 8.1.8.---}\enspace
\emph{Let \(S\) and \(S'\) be affine schemes, let \(S'\to S\) be
faithfully flat, let \(J\) be a square-zero ideal on \(S\), and let
\(S_0\) be the closed subscheme which it defines. Let \(G\) be a
smooth \(S\)-group. Put
\[
  S'_0=S'\times_SS_0,
  \qquad
  G_0=G\times_SS_0.
\]
Then the canonical map
\[
  H^1(S'/S,G)\longrightarrow H^1(S'_0/S_0,G_0)
\]
is bijective.}

\par\medskip\noindent\SGAInlineRefTarget{sga3:xxiv:remark:8-1-9}\textbf{Remark 8.1.9.---}\enspace
If \(G\) is commutative, the same assertion holds for every
\(H^i\), \(i>0\), with an analogous proof.

\smallskip\noindent\emph{Proof.}
Let \(M_0\) be the quasi-coherent \(\mathcal O_{S_0}\)-module
\[
  M_0
  =
  \Lie(G_0/S_0)\otimes_{\mathcal O_{S_0}}J.
\]
For every \(S_0\)-scheme \(T_0\), put
\[
  M_0(T_0)
  =
  H^0\!\left(
    T_0,
    M_0\otimes_{\mathcal O_{S_0}}\mathcal O_{T_0}
  \right),
\]
and define the \(S\)-group functors
\[
  M=\prod_{S_0/S}M_0,
  \qquad
  \overline G=\prod_{S_0/S}G_0
\]
by
\[
  M(T)=M_0(T_0),
  \qquad
  \overline G(T)=G_0(T_0),
  \qquad
  T_0=T\times_SS_0.
\]
By \SGARefLink{sga3:iii:corollary:0-9}{Exposé III, 0.9} and (\SGARefLink{sga3:iii:remark:0-6:equation:0-6-2}{0.6.2}), for every affine \(S\)-scheme \(T\)
there is a functorial exact sequence
\[
  1\longrightarrow M(T)\longrightarrow G(T)
  \longrightarrow\overline G(T)\longrightarrow1.
\]
We must study the map
\[
  H^1(S'/S,G)
  \longrightarrow
  H^1(S'_0/S_0,G_0)
  =
  H^1(S'/S,\overline G).
\]

Suppose first that \(G\) is commutative. There is then an exact
cohomology sequence
\[
  H^1(S'/S,M)
  \longrightarrow H^1(S'/S,G)
  \longrightarrow H^1(S'/S,\overline G)
  \longrightarrow H^2(S'/S,M).
\]
Moreover,
\[
  H^i(S'/S,M)
  =
  H^i(S'_0/S_0,M_0)
  =
  H^i(S'_0/S_0,\mathcal M_0)
  =
  0
  \qquad(i\neq0),
\]
the last equality following from TDTE I, B, Lemma~1.1.

Now suppose that \(G\) is not commutative. We use the exact sequence
of nonabelian cohomology. If
\(u\in Z^1(S'/S,\overline G)\), the elements of \(H^1(S'/S,G)\) whose
image in \(H^1(S'/S,\overline G)\) equals the class of \(u\) lie in the
image of the corresponding coboundary map
\[
  H^1(S'/S,M_u)\longrightarrow H^1(S'/S,G),
\]
where \(M_u\) is \(M\) twisted by \(u\). Similarly, for
\(v\in Z^1(S'/S,\overline G)\), there is an obstruction
\[
  \Delta(v)\in H^2(S'/S,M_v)
\]
which vanishes exactly when the class of \(v\) is in the image of
\(H^1(S'/S,G)\). It therefore suffices to prove
\[
  H^1(S'/S,M_u)=0,
  \qquad
  H^2(S'/S,M_v)=0.
\]

Recall (\SGARefLink{sga3:iii:proposition:0-8}{Exposé III, 0.8}) that the action of \(G\) on \(M\) defined by
the exact sequence above is induced functorially by the adjoint
representation
\[
  \operatorname{ad}:G_0
  \longrightarrow
  \underline{\operatorname{Aut}}_{\mathcal O_{S_0}}
    (\Lie(G_0/S_0)).
\]
The cocycle \(u\), respectively \(v\), therefore defines a descent
datum on \(\Lie(G_0/S_0)\); write \(L_u\), respectively \(L_v\), for
the resulting quasi-coherent \(\mathcal O_{S_0}\)-module. For
\(T\to S\),
\[
  M_u(T)
  =
  H^0\!\left(
    T_0,
    L_u\otimes_{\mathcal O_{S_0}}J
      \otimes_{\mathcal O_{S_0}}\mathcal O_{T_0}
  \right),
\]
and similarly with \(v\) in place of \(u\). Hence
\[
  H^1(S'/S,M_u)=H^1(S'_0/S_0,L_u\otimes J),
\]
\[
  H^2(S'/S,M_v)=H^2(S'_0/S_0,L_v\otimes J),
\]
and both groups vanish by the result already used.
\hfill\(\square\)

\par\medskip\noindent\SGAInlineRefTarget{sga3:xxiv:proposition:8-2}\textbf{Proposition 8.2.---}\enspace
\emph{Let \(\mathcal C\) be a category with fiber products, endowed
with a topology coarser than the canonical topology. Let \(S'\to S\)
be a morphism of \(\mathcal C\), let \(G'\) be an \(S'\)-group sheaf,
and let
\[
  G=\prod_{S'/S}G'
\]
be the corresponding \(S\)-group sheaf. Let
\[
  H^1_S(S',G')\subset H^1(S',G')
\]
be the subset of classes of principal homogeneous sheaves under
\(G'\) which are trivialized by a sieve on \(S'\) obtained by base
change from a suitable covering sieve on \(S\). The canonical map
\[
  H^1(S,G)\longrightarrow H^1(S',G'),
  \qquad
  P\longmapsto P\times_SS',
\]
induces a bijection
\[
  H^1(S,G)\xrightarrow{\sim}H^1_S(S',G').
\]
Its inverse is induced by
\[
  P'\longmapsto\prod_{S'/S}P'.
\]}

For every object \(X\) of \(\mathcal C/S\), the definition gives an
isomorphism, functorial in \(X\),
\[
  \operatorname{Hom}_S(X,G)
  \xrightarrow{\sim}
  \operatorname{Hom}_{S'}(X\times_SS',G').
\]
Thus, for every \(S\)-object \(T\), there is a bijection functorial in
\(T\),
\[
  H^1(T/S,G)
  \xrightarrow{\sim}
  H^1(T'/S',G'),
  \qquad
  T'=T\times_SS'.
\]
Replacing the single morphism \(T\to S\) by an arbitrary covering
family of \(S\), and then passing to the inductive limit, gives the
first part of the proposition. The second follows immediately.

\par\medskip\noindent\SGAInlineRefTarget{sga3:xxiv:lemma:8-3}\textbf{Lemma 8.3.---}\enspace
\emph{Under the hypotheses of \SGARefLink{sga3:xxiv:proposition:8-2}{8.2}, the equality
\[
  H^1_S(S',G')=H^1(S',G')
\]
is local on \(S\). More precisely, suppose that there is a covering
family \(\{S_i\to S\}\) such that, for every \(i\),
\[
  H^1_{S_i}(S'\times_SS_i,G')
  =
  H^1(S'\times_SS_i,G').
\]
Then \(H^1_S(S',G')=H^1(S',G')\).}

Indeed, let \(P'\) be a principal homogeneous sheaf under \(G'\), and
put
\[
  P'_i=P'\times_{S'}(S'\times_SS_i).
\]
By hypothesis, there is a covering family
\(\{S_{ij}\to S_i\}\) such that, for each \(j\),
\[
  P'\times_{S'}(S'\times_SS_{ij})
\]
has a section. But \(\{S_{ij}\to S\}\) is a covering family of \(S\),
and \(P'\) is therefore trivialized by the covering family of \(S'\)
obtained from it by base change.

\par\medskip\noindent\SGAInlineRefTarget{sga3:xxiv:proposition:8-4}\textbf{Proposition 8.4.---}\enspace
\emph{Let \(S'\to S\) be a finite étale morphism of schemes. Let
\(G'\) be an \(S'\)-group sheaf and let
\[
  G=\prod_{S'/S}G'
\]
be the corresponding \(S\)-group sheaf. For the étale topology,
respectively the local finite étale topology, respectively the
\textup{(fpqc)} topology, the following functors induce mutually
inverse bijections.}

\[
  \begin{tikzcd}[column sep=huge,row sep=large]
\mathcal P
  \arrow[r,mapsto]
&
\mathcal P\times_S S'
\\
\displaystyle\prod_{S'/S}\mathcal P'
&
\mathcal P'
  \arrow[l,mapsto]
\end{tikzcd}

\]

\emph{In particular,
\[
  H^1(S,G)\simeq H^1(S',G').
\]}

By \SGARefLink{sga3:xxiv:proposition:8-2}{8.2}, it is enough to show
\[
  H^1_S(S',G')=H^1(S',G').
\]
By \SGARefLink{sga3:xxiv:lemma:8-3}{8.3}, this may be checked locally for the local finite étale
topology. We may therefore suppose that \(S'\) is a finite disjoint
union of copies of \(S\), say \(I_S\), where \(I\) is a finite set.
Then \(G'\) is given by a family \((G_i)_{i\in I}\) of sheaves on
\(S\), and
\[
  H^1(S',G')
  \simeq
  \prod_{i\in I}H^1(S,G_i).
\]
On the other hand,
\[
  H^1(S,G)
  \simeq
  \prod_{i\in I}H^1(S,G_i).
\]
It follows from \SGARefLink{sga3:xxiv:proposition:8-2}{8.2} that
\(H^1_S(S',G')=H^1(S',G')\).
\hfill\(\square\)

\par\medskip\noindent\SGAInlineRefTarget{sga3:xxiv:remark:8-5}\textbf{Remarks 8.5.---}\enspace
\SGARefLink{sga3:xxiv:proposition:8-2}{Propositions 8.2} and \SGARefLink{sga3:xxiv:proposition:8-4}{8.4} may be interpreted through the following
exact sequence, where \(f:S'\to S\) is the given morphism:
\[
  1\longrightarrow H^1(S,f_*G')
  \longrightarrow H^1(S',G')
  \longrightarrow H^0(S,R^1f_*G').
\]
In the commutative case, this exact sequence follows from the Leray
spectral sequence. It remains valid in the noncommutative
case.\footnote{\textit{Editors' note.} See \S V.3.1.4 of
\([{\rm Gi71}]\).}

In this form one sees that the result remains valid when \(f\) is
assumed only finite, or even merely integral, provided that the
topology is the étale topology. Indeed, for every \(G'\),
\[
  R^1f_*(G')=\text{the final sheaf}
\]
by SGA~4, VIII, 5.3.\footnote{\textit{Editors' note.} This reference
also points to \S V.3.1.4 of \([{\rm Gi71}]\).}

The result becomes false for a topology such as \textup{(fpqc)} or
\textup{(fppf)}, even if
\[
  S=\operatorname{Spec}(k),
  \qquad
  S'=\operatorname{Spec}(k[t]/t^2),
\]
where \(k\) is an algebraically closed field of characteristic
\(p\neq0\), and \(G'=\mu_p\) or \(\alpha_p\).

Likewise, 8.2 becomes false, even for the étale topology, if the
finiteness hypothesis on \(f\) is omitted. For example, take \(f\)
to be an open immersion with
\[
  S=\operatorname{Spec}(V),
\]
where \(V\) is a complete discrete valuation ring with algebraically
closed residue field, take \(S'\) to be the open subscheme consisting
of the generic point, and take \(G'\) to be the constant group
\((\mathbb Z/n\mathbb Z)_{S'}\), where \(n\) is prime to the residue
characteristic of \(V\). Then
\[
  H^1(S,G)=0,
  \qquad
  H^1(S',G')\neq0.
\]
Replacing \(S'\) by \(S\amalg S'\) gives an analogous example in which
\(S'\to S\) is étale and surjective, hence a cover for the topology
under consideration.

\section*{Bibliography}

\noindent[BLie]\enspace
N.~Bourbaki, \textit{Groupes et algèbres de Lie}, Chap.~I, Hermann,
1960.\par

\smallskip\noindent[Ch51]\enspace
C.~Chevalley, \textit{Théorie des groupes de Lie}, t.~II,
\textit{Groupes algébriques}, Hermann, 1951.\par

\smallskip\noindent[Di57]\enspace
J.~Dieudonné, Lie groups and Lie hyperalgebras over a field of
characteristic \(p>0\). VI, \textit{Amer. J. Math.} \textbf{79}
(1957), no.~2, 331--388.\footnote{\textit{Editors' note.} The editors
have added, after these three references which appeared in the
original, the references that follow.}\par

\smallskip\noindent[Bo91]\enspace
A.~Borel, \textit{Linear Algebraic Groups}, 2nd ed.,
Springer-Verlag, 1991.\par

\smallskip\noindent[Ch05]\enspace
C.~Chevalley, \textit{Classification des groupes algébriques
semi-simples}, with the collaboration of P.~Cartier,
A.~Grothendieck, and M.~Lazard, \textit{Collected Works}, vol.~3,
Springer, 2005.\par

\smallskip\noindent[DG70]\enspace
M.~Demazure and P.~Gabriel, \textit{Groupes algébriques}, Masson and
North-Holland, 1970.\par

\smallskip\noindent[Gi71]\enspace
J.~Giraud, \textit{Cohomologie non abélienne}, Springer-Verlag,
1971.\par

\smallskip\noindent[Ja87]\enspace
J.~C.~Jantzen, \textit{Representations of Algebraic Groups}, Academic
Press, 1987; 2nd ed., American Mathematical Society, 2003.\par

\smallskip\noindent[Jou83]\enspace
J.-P.~Jouanolou, \textit{Théorèmes de Bertini et applications},
Birkhäuser, 1983.\par

\smallskip\noindent[Ta75]\enspace
M.~Takeuchi, On coverings and hyperalgebras of affine algebraic
groups, \textit{Trans. Amer. Math. Soc.} \textbf{211} (1975),
249--275.\par

